\scanpage{163}
\mspage{18}
où $\widehat{\widetilde F}$ est le \uline{complété formel} de $\widetilde F$ le long de $\widetilde X-X$ (nb fini de points\ldots). C'est ici, sauf erreur, que \unsure{s'amorce|s'annonce} le lien avec le corps de classes local, version cohomologique, dont je vais dire quelques mots maintenant.

E) \uline{Corps de classes local} comme formule de dualité.

Soit $V$ un anneau de val. discrète complet, à corps résiduel $k$ -- on suppose ou bien $k$ relevé en $k\subset V$ (\unsure{donc} $V\simeq\unsure{h|k}[[T]]$) ou bien $k$ parfait. Pour fixer les idées et être sur un terrain solide, je prends \inserted{d'abord} sur $K={}$corps des fractions de $V$ des coefficients \uline{finis} $F$ (comme sur $X$ tout à l'heure), et je m'intéresse aux $H^i(K,F)$ ou encore à $R\Gamma_K(F)$. L'essentiel du travail à faire consiste à définir une catégorie de coeff. convenable sur $k$ (sans doute la $\overline{\mathrm m}$ que dans D)) et un foncteur
\[
 F\longmapsto\mathbb{R}\underline{\Gamma}_K(F)\unsure{:|.}
\]
vers ces coefficients, de telle façon que l'on ait
\[
 R\Gamma_K(F)\simeq R\Gamma_h\bigl(R\underline{\Gamma}_K(F)\bigr).
\]
% archive sheet number: 163
